\subsubsection{DHAT}\label{valgrind-dhat}

\href{https://valgrind.org/docs/manual/dh-manual.html}{DHAT} je Valgrind alat za analizu ponašanja hip alokacija. Koristan je kada program nema očigledno curenje memorije, ali pravi previše kratkoživećih alokacija, često realocira iste podatke ili koristi mali deo alociranog bloka.

DHAT ne zamenjuje Memcheck. Memcheck prvo treba da potvrdi da program nema neispravna čitanja, pisanja i curenja. Nakon toga DHAT prikazuje kako se memorija koristi: koliko blokova je alocirano, koliko bajtova prolazi kroz alokator, koliko dugo blokovi žive i koliko se često čitaju ili pišu.

Primer u datoteci \texttt{dhat\_churn\_bad.c} pravi veliki broj malih privremenih stringova. Svaki prolaz alocira niz pokazivača, zatim za svaki zapis alocira zaseban string i odmah ga oslobađa na kraju prolaza.

\codefile{dhat_churn_bad.c}{c}

Popravljena verzija zadržava isti izlazni radni tok, ali koristi jedan radni bafer za sve zapise.

\codefile{dhat_churn_fixed.c}{c}

Programi se prevode i pokreću ovako:

\begin{minted}{bash}
$ make
$ valgrind --tool=dhat ./dhat_churn_bad.out
$ valgrind --tool=dhat ./dhat_churn_fixed.out
\end{minted}

DHAT pravi datoteku oblika \texttt{dhat.out.\textless{}pid\textgreater{}}. Izveštaj se otvara kroz HTML prikazivač koji dolazi uz Valgrind:

\begin{minted}{bash}
$ firefox /usr/lib/valgrind/dh_view.html
\end{minted}

Putanja do \texttt{dh\_view.html} zavisi od distribucije. Ako komanda iznad ne postoji, pronađite je komandom:

\begin{minted}{bash}
$ find /usr -name dh_view.html 2>/dev/null
\end{minted}

U poređenju ova dva programa obratiti pažnju na sledeće stavke u DHAT izveštaju:

\begin{itemize}
\tightlist
\item
  \textbf{Total blocks}: loša verzija pravi mnogo više alokacija.
\item
  \textbf{Total bytes}: loša verzija propušta više bajtova kroz alokator.
\item
  \textbf{At t-gmax}: pokazuje mesto najveće zauzetosti hipa.
\item
  \textbf{Reads/Writes}: pomaže da se vidi da li je alocirani blok zaista korišćen ili je samo rezervisan.
\end{itemize}

Za vežbu prvo proveriti obe verzije Memcheck-om, zatim uporediti DHAT izveštaje. Očekivani zaključak je da popravljena verzija ne menja semantiku programa, ali značajno smanjuje broj hip alokacija i ukupan promet kroz alokator.
